\section{Introduction}\label{sec:introduction}

This chapter surveys existing approaches to music representation and programming, focusing on systems that enable the
creation, manipulation, and execution of musical structures through textual or programmatic means.

The reviewed works span multiple categories, including high-level composition systems, low-level audio programming
languages, and music representation formats.
These systems differ significantly in their design goals, levels of abstraction, and approaches to modeling musical
time and structure.

High-level composition systems, such as \textbf{Sonic Pi}, \textbf{Tidal}, and \textbf{HMusic}, emphasize expressive
musical abstractions, often prioritizing accessibility, structured composition, and interactive use.

In contrast, low-level audio programming languages like \textbf{ChucK} and \textbf{SuperCollider} provide fine-grained
control over sound synthesis and timing, but operate at a lower level of abstraction.

Additionally, music representation formats such as \textbf{MIDI} and \textbf{ABC notation} define how musical data can
be encoded, stored, and transmitted, often prioritizing interoperability and readability over compositional abstraction.

By analyzing these systems, this chapter identifies key design trade-offs and limitations in existing approaches,
particularly with respect to temporal modeling, abstraction level, and compositional structure.
These observations serve as the foundation for motivating the design of the domain-specific language proposed in
this thesis.
